% Part: incompleteness
% Chapter: introduction
% Section: undecidability

\documentclass[../../../include/open-logic-section]{subfiles}

\begin{document}

\olfileid[id]{inc}{int}{dec}

\olsection{Ketakterputusan dan Ketaklengkapan}

Pembuktian G\"odel atas teorema-teorema ketaklengkapan memerlukan
aritmetisasi sintaksis. Akan tetapi, bahkan tanpa aritmetisasi itu, kita dapat
memperoleh beberapa hasil yang baik hanya dengan asumsi bahwa suatu teori
!!{represents} semua relasi !!{decidable}. Pembuktiannya merupakan argumen
diagonal yang serupa dengan pembuktian ketakterputusan masalah penghentian.

\begin{thm}
Jika $\Gamma$ adalah teori konsisten yang !!{represents} setiap relasi
!!{decidable}, maka $\Gamma$ tidak !!{decidable}.
\end{thm}

\begin{proof}
Andaikan $\Gamma$ !!{decidable}. Kita tunjukkan bahwa jika $\Gamma$
!!{represents} setiap relasi !!{decidable}, teori itu pasti tak konsisten.

Sifat !!^{decidable} (relasi satu-tempat) direpresentasikan oleh
!!{formula}s dengan satu variabel bebas. Misalkan $!A_0(x)$, $!A_1(x)$,
\dots, merupakan enumerasi yang dapat dihitung atas semua !!{formula}s
semacam itu. Sekarang perhatikan himpunan $D \subseteq \Nat$ berikut:
\[
D = \Setabs{n}{\Gamma \Proves \lnot !A_n(\num{n})}
\]
Himpunan $D$ !!{decidable}, sebab kita dapat menguji apakah $n \in D$
dengan mula-mula menghitung $!A_n(x)$, lalu dari situ menghitung $\lnot
!A_n(\num{n})$. Jelas bahwa menyubstitusikan suku $\num{n}$ bagi setiap
kemunculan bebas $x$ dalam $!A_n(x)$ dan mengawali $!A_n(\num{n})$ dengan
$\lnot$ merupakan prosedur mekanis. Berdasarkan asumsi, $\Gamma$
!!{decidable}, sehingga kita dapat menguji apakah $\lnot !A_n(\num{n}) \in
\Gamma$. Jika ya, $n \in D$; jika tidak, $n \notin D$. Jadi $D$ juga
!!{decidable}.

Karena $\Gamma$ !!{represents} semua sifat !!{decidable}, teori itu
!!{represents}~$D$. Semua !!{formula}s yang !!{represents}s $D$
dalam~$\Gamma$ terdapat di antara $!A_0(x)$, $!A_1(x)$, \dots. Maka,
misalkan $d$ adalah bilangan sedemikian sehingga $!A_d(x)$ !!{represents}
$D$ dalam~$\Gamma$. Jika $d \notin D$, karena $!A_d(x)$ !!{represents}~$D$,
$\Gamma \Proves \lnot !A_d(\num{d})$. Namun, ini berarti bahwa $d$
memenuhi syarat penentu~$D$, sehingga $d \in D$. Hal ini bertentangan dengan
$d \notin D$. Jadi, melalui pembuktian tak langsung, $d \in D$.

Karena $d \in D$, menurut definisi~$D$, $\Gamma \Proves \lnot
!A_d(\num{d})$. Di pihak lain, karena $!A_d(x)$ !!{represents}~$D$
dalam $\Gamma$, $\Gamma \Proves !A_d(\num{d})$. Dengan demikian,
$\Gamma$ tak konsisten.
\end{proof}

\begin{explain}
Teorema sebelumnya menunjukkan bahwa tidak ada teori konsisten yang
!!{represents} semua relasi !!{decidable} yang sekaligus !!{decidable}.
Kita akan menunjukkan bahwa $\Th{Q}$ memang !!{represents}s semua relasi
!!{decidable}; artinya, semua teori yang memuat $\Th{Q}$, seperti $\Th{PA}$
dan $\Th{TA}$, juga demikian dan karenanya juga tidak !!{decidable}.
(Karena semua teori ini benar dalam model standar, semuanya konsisten.)

Hasil ini juga dapat kita gunakan untuk memperoleh versi lemah teorema
ketaklengkapan pertama. Setiap teori yang !!{axiomatizable} dan !!{complete}
adalah !!{decidable}. Karena itu, teori konsisten yang !!{axiomatizable} dan
!!{represents}s semua sifat !!{decidable} tidak mungkin !!{complete}.
\end{explain}

\begin{thm}
Jika $\Gamma$ !!{axiomatizable} dan !!{complete}, maka teori itu
!!{decidable}.
\end{thm}

\begin{proof}
Setiap teori tak konsisten adalah !!{decidable}, sebab teori tak konsisten
memuat semua !!{sentence}s; jadi jawaban atas pertanyaan ``apakah $!A \in
\Gamma$'' selalu ``ya,'' sehingga dapat diputuskan.

Sekarang andaikan $\Gamma$ konsisten, serta !!{axiomatizable} dan
!!{complete}. Karena $\Gamma$ !!{axiomatizable}, teori itu
!!{computably enumerable}. Kita dapat mengenumerasi semua !!{derivation}s
yang benar dari aksioma-aksioma~$\Gamma$ dengan suatu fungsi yang dapat
dihitung. Dari suatu !!{derivation} yang benar, kita dapat menghitung
!!{sentence} hasilnya karena derivasi tersebut !!{derive}s kalimat itu;
dengan menggabungkan keduanya, diperoleh fungsi yang
dapat dihitung yang mengenumerasi semua teorema~$\Gamma$. Untuk suatu
!!{sentence}~$!A$, kalimat tersebut merupakan teorema~$\Gamma$ jika dan
hanya jika $\lnot !A$ bukan teorema,
sebab $\Gamma$ konsisten dan !!{complete}. Karena itu, kita dapat memutuskan
apakah $!A \in \Gamma$ sebagai berikut. Enumerasikan semua teorema
$\Gamma$. Ketika $!A$ muncul dalam daftar, kita mengetahui bahwa $\Gamma
\Proves !A$. Ketika $\lnot !A$ muncul dalam daftar, kita mengetahui bahwa
$\Gamma \Proves/ !A$. Karena $\Gamma$ !!{complete}, salah satu kasus ini
pada akhirnya terjadi, sehingga prosedur tersebut akhirnya menghasilkan
jawaban.
\end{proof}

\begin{cor}
\ollabel{cor:incompleteness}
Jika $\Gamma$ konsisten, !!{axiomatizable}, dan !!{represents} setiap sifat
!!{decidable}, maka teori itu tidak !!{complete}.
\end{cor}

\begin{proof}
Jika $\Gamma$ !!{complete}, teori itu akan !!{decidable} menurut teorema
sebelumnya (karena teori itu !!{axiomatizable} dan konsisten). Namun, karena
$\Gamma$ !!{represents} setiap sifat !!{decidable}, teori itu tidak
!!{decidable} menurut teorema pertama.
\end{proof}

\begin{prob}
Tunjukkan bahwa $\Th{TA} = \Setabs{!A}{\Sat{N}{!A}}$ tidak
!!{axiomatizable}. Anda boleh mengasumsikan bahwa $\Th{TA}$ merepresentasikan
semua sifat yang dapat diputuskan.
\end{prob}

Setelah kita menetapkan bahwa, misalnya, $\Th{Q}$ !!{represents} semua sifat
!!{decidable}, korolari tersebut memberi tahu kita bahwa $\Th{Q}$ pasti tak
lengkap. Namun, pembuktiannya tidak memberikan contoh !!{sentence}
independen; pembuktian itu hanya menunjukkan bahwa !!a{sentence} semacam itu
pasti ada. Untuk memperoleh contohnya, kita harus mengaritmetisasi sintaksis
dan mengikuti gagasan pembuktian asli G\"odel. Tentu saja, kita juga masih
harus menunjukkan klaim pertama, yakni bahwa $\Th{Q}$ memang
!!{represents}s semua sifat !!{decidable}.

Perlu dicatat bahwa tidak setiap teori yang \emph{menarik} bersifat tak
lengkap atau tak dapat diputuskan. Ada banyak teori yang cukup kuat untuk
mendeskripsikan fakta matematis yang menarik, tetapi tidak memenuhi syarat
hasil G\"odel. Sebagai contoh, $\Th{Pres} = \Setabs{!A \in
\Lang{L_{A^+}}}{\Sat{N}{!A}}$, yakni himpunan !!{sentence}s dalam bahasa
aritmetika tanpa~$\times$ yang benar dalam model standar, bersifat lengkap
dan dapat diputuskan. Teori ini disebut aritmetika Presburger dan membuktikan
semua kebenaran tentang bilangan asli yang dapat dirumuskan hanya dengan
$\Obj{0}$, $\prime$, dan~$+$.

\end{document}
